\section{Обсуждение}
\label{sec:discussion}

Иерархия факторов точности симулятора такова: (1)~тип операции~--- GEMM против
нормализации~--- определяет четыре порядка разрыва; (2)~внутри GEMM доминирует
компилятор, а не реализация; (3)~тип ошибки ($b$) классифицируется за один
проход; (4)~накопление по слоям определяется архитектурой (суммирующая связь,
тип нормализации, насыщающиеся нелинейности) сильнее, чем типом ошибки.

Ряд результатов контринтуитивен. Относительное расхождение GEMM почти не растёт с
размерностью редукции $K$ (отношение $1.05$ вместо $\sqrt{3.5}\approx1.87$),
что согласуется с доминированием компилятора, а не GEMM-арифметики. Данное
наблюдение получено на двух значениях $K$ и не является формальным опровержением
степенной зависимости~\eqref{eq:scaling}; оно служит косвенным указанием,
проверяемым в разделе~\ref{sec:results} независимыми средствами (разделение
вкладов компилятора и GEMM). cuBLAS оказался ближе к CPU, чем CUDA-ядро, несмотря
на явный контроль порядка накопления в последнем. В изолированном режиме обратный
проход накапливает сопоставимо с прямым (предсказание $\text{BWD}>\text{FWD}$ не
подтвердилось); исключение~--- DeepNorm
($R_{\text{bwd}}=0.35\ll R_{\text{fwd}}=1.39$).

\textbf{Где повышать точность.} Усиление расхождения слоем ограничено числом
обусловленности его весовой матрицы $\kappa(W)$~--- отношением наибольшего и
наименьшего сингулярных чисел (раздел~\ref{sec:model}). Это даёт прикладной
критерий: наибольший эффект от повышения точности вычислений дают слои с большим
$\kappa(W)$, не имеющие сжимающего стока на выходе; слои за нормализацией,
насыщающейся нелинейностью или усреднением гасят расхождение независимо от
реализации, и затраты на их уточнение неоправданны. В обследованных компонентах
обученные веса дают умеренное $\kappa(W)$, что и объясняет ограниченное
накопление и обосновывает применимость линейного приближения.

\textbf{Следствия для обучения.} Результаты обратного прохода переносятся на
обучение через шаг оптимизатора $\Vec{w}\leftarrow\Vec{w}-\eta(\Vec{g}+\delta)$,
где $\delta$~--- расхождение градиента. За $T$ шагов смещение параметров
накапливается как ${\sim}\,T\eta\|\delta\|$, поэтому решающее значение имеет не
норма, а наличие систематической составляющей: случайная ($b\approx0.5$) ошибка
усредняется по мини-батчам и шагам и действует как дополнительный шум градиента,
к которому обучение устойчиво, тогда как смещённая ($b$ вдали от $0.5$, как у
cuBLAS) направленно сдвигает оценку первого момента у адаптивных оптимизаторов и
траекторию обучения. Это согласуется с первичностью \emph{типа} ошибки над её
величиной и обосновывает предпочтение реализаций без смещения (CPU-эталон,
CUDA-ядро) в конфигурациях обучения.

\textbf{Ограничения.} Эксперименты выполнены в режиме одного запуска: значения
$d_{\text{rel}}^{(l)}$ и $R$ воспроизводятся побитово, но без оценки дисперсии по
ансамблю входов; малые межкомпонентные различия агрегированных метрик ($R$, $b$,
$c$) трактуются как наблюдательные. Естественное продолжение~--- статистический
режим с усреднением по ансамблю входов и инициализаций, дающий доверительные
интервалы для $R$, $b$, $c$. Рассмотрены 2--6-слойные блоки, а не полноглубинные
модели; авторегрессивные и рекуррентные архитектуры вне работы. Все вычисления
велись в \texttt{FP32}; в режимах пониженной разрядности (\texttt{TF32},
\texttt{BF16}) арифметический вклад реализации растёт и соотношение с
компиляторным может измениться~--- это отдельное направление. Наконец, близость к
CPU-эталону не равнозначна близости к арифметике конкретного устройства~---
последнее требует верификации на реальном оборудовании.
